\section{Coefficients, stabilizers, and comparison inputs}\label{h3:sec:framework}

The coefficient arguments require three kinds of data: the local sphere
and its finite constant-field extensions, the stabilizer actions on the
deformation rings, and the topology used on continuous cochains.  We fix
these data here and state the external bundle and boundary results at the
precise scope used later.  The actual comparisons between algebraic and
completed coefficients will be proved in the following sections.

\subsection{Localizations and coefficient fields}

Fix \(p\geq5\). We work with \(p\)-local spectra and write
\[
 S=\mathbb S_p^\wedge,\qquad T=L_{K(3)}S,\qquad
 X=L_2T,\qquad R_i=L_{K(i)}S.
\]
We use the natural \(S\)-module structures throughout. In particular,
\(L_0S=H\Q_p\), and the rational maps used in fracture are
\(H\Q_p\)-linear. The rationalization of \(L_iS\), for \(i=1,2\), is
\(H\Q_p\); it should be distinguished from \(L_0R_i\).

Choose a finite field \(k\supseteq\F_{p^6}\). Further finite enlargement
will be specified when needed to define markings or characters.
Let \(W(k)\) be its Witt ring. The finite étale extension
\(\Z_p\to W(k)\) has a unique finite étale lift to an \(E_\infty\)
\(S\)-algebra \(S_k\), with
\[
 \pi_*S_k=\pi_*S\otimes_{\Z_p}W(k).
\]
This is the finite étale classification for ring spectra;
see \cite[Theorem~2.32]{Mathew}. We write \(R_{i,k}=L_{K(i)}S_k\).
A \(\Z_p\)-basis of \(W(k)\) gives an equivalence of the underlying
\(S\)-module with a finite wedge of copies of \(S\), so extension to
\(S_k\) is faithful. The Galois automorphisms lift coherently through
the same classification. They will be used to descend constructions
made over \(k\).

Let \(E_n(k)\) be Morava \(E\)-theory for the height-\(n\) Honda formal
group over \(k\). At height three our coefficient convention is
\[
 E_{3,*}(k)=W(k)[[x,y]][u^{\pm1}],\qquad |u|=-2.
\]
The invariant cotangent, or periodicity, line is
\(\omega=\pi_2E_3(k)\); it is generated by \(u^{-1}\) over the
degree-zero ring. Set
\[
 A=k[[x,y]],\qquad x=u_1,\quad y=u_2,\qquad B_0=A[x^{-1}].
\]
Twists by powers of \(\omega\) and finite characters are retained
whenever a finiteness argument requires them. A trivialization made
after finite extension will always be descended.

\subsection{Stabilizers and reduced norm}

Let \(D_n\) be the central division algebra over \(\Q_p\) of invariant
\(1/n\), and let
\[
 P_n=\mathcal O_{D_n}^{\times}.
\]
This is the ordinary Honda stabilizer. Its action on \(E_n(k)\) fixes
\(W(k)\). The coefficient-field Galois group acts semilinearly and gives
the extended stabilizer
\[
 G_n(k)=P_n\rtimes\Gal(k/\F_p).
\]
The field contains the endomorphism field for every \(n=1,2,3\) used here.

At height three, define
\[
 H=\ker\bigl(\hthreeNrd:P_3\to\Z_p^\times\bigr),\qquad
 H'=\hthreeNrd^{-1}(\mu_{p-1}).
\]
The reduced norm is surjective, and therefore
\begin{equation}\label{h3:eq:framework-norm-quotients}
 H'/H\simeq\mu_{p-1},\qquad
 P_3/H'\simeq1+p\Z_p\simeq\Z_p.
\end{equation}
For example, the norm on the maximal unramified subfield is surjective
on units. The logarithm on the last quotient is normalized using a
topological generator. The resulting integral determinant class is
denoted \(\zeta\in\pi_{-1}T\); its construction as an actual sphere map
is recalled in \Cref{h3:sec:descent}.

\begin{lemma}\label{h3:lem:framework-groups}
The compact groups \(P_3,H,P_2,P_1,\GL_2(\Z_p)\) and
\(\hthreeSL_2(\Z_p)\) have no elements of order \(p\). Their \(p\)-cohomological
dimensions are respectively \(9,8,4,1,4,3\). Their adjoint orientation
characters are trivial. Their trivial \(\Z_p\)-modules admit finite
resolutions by finitely generated projective completed group-ring
modules.
\end{lemma}

\begin{proof}
An element of order \(p\) in \(D_n^\times\) would generate a copy of
\(\Q_p(\zeta_p)\). Its degree \(p-1\) must divide \(n^2\), since
\(D_n\) would be a vector space over that subfield. For \(n=1,2,3\)
and \(p\geq5\), this is impossible except for the apparent numerical
possibility \(n=2,p=5\); in that case a commutative subfield of a
degree-two central division algebra has degree at most two, whereas
\(p-1=4\). Equivalently, the usual subfield-degree theorem rules out
all these cases at once. An element of order \(p\) in
\(\GL_2(\Q_p)\) would give a faithful two-dimensional representation
of the nontrivial irreducible factor of \(X^p-1\), whose degree is
\(p-1>2\), and is likewise impossible.

The Lie dimensions are the stated ones. Reduced norm has surjective
differential, so its kernel has dimension eight. Conjugation has
determinant one on a central simple algebra; its restriction to the
trace-zero summand also has determinant one. The same calculation
applies to the matrix groups. Thus their adjoint orientation
characters are trivial.

For a compact \(p\)-adic analytic group without \(p\)-torsion,
\(p\)-cohomological dimension equals its Lie dimension
\cite[Section~1, Corollary~(1)]{Serre1965CohomologicalDimension}.
Writing \(\Lambda=\Z_p[[G]]\), Venjakob
\cite[pp.~275--276]{Venjakob} records both Noetherianity of \(\Lambda\)
and \(\operatorname{pd}_{\Lambda}\Z_p=\operatorname{cd}_p(G)\),
with agreement between finitely generated compact and abstract
modules.  Successive finite free presentations therefore have finitely
generated kernels and a projective final syzygy, giving the asserted
finite projective resolution. The associated duality is used with the
continuous coefficient models specified below.
\end{proof}

No torsion-freeness claim about an arbitrary finite Galois factor is
needed. Its action is instead handled by semilinear descent.

\begin{lemma}\label{h3:lem:framework-semilinear}
Let \(k'/k\) be a finite extension of finite fields with Galois group
\(\Gamma\). On \(W(k')\)-modules carrying a semilinear \(\Gamma\)-action,
the invariants functor is exact. The same assertion holds for reductions
modulo \(p^\ell\), and the associated finite-group higher cohomology
vanishes. These assertions respect continuous equivariant maps.
\end{lemma}

\begin{proof}
The residue-field trace is surjective. Lifting a trace-one element to
\(W(k')\) and dividing by its unit trace produces \(a\in W(k')\) with
\(\sum_{\gamma\in\Gamma}\gamma(a)=1\). For a semilinear module \(M\), the
finite sum
\[
 P(m)=\sum_{\gamma\in\Gamma}\gamma(a)\,\gamma(m)
\]
is an additive projection onto \(M^\Gamma\). If an invariant element
in a quotient has a lift, applying \(P\) to that lift gives an invariant
lift. This proves exactness, and the same formula works modulo
\(p^\ell\). For an equivariant homogeneous cochain \(f\) of positive
degree, the operator
\[
 (hf)(g_0,\ldots,g_{q-1})
 =\sum_{\gamma\in\Gamma}\gamma(a)
 f(\gamma,g_0,\ldots,g_{q-1})
\]
is still equivariant, by semilinearity and reindexing the sum. The
alternating homogeneous differential satisfies \(dh+hd=\hthreeid\) in
positive degrees, since the sum of the coefficients is one. This proves
the higher vanishing. All sums and scalar operations are finite and
continuous. No division by \(|\Gamma|\) has been made.
\end{proof}

\subsection{Continuous cochains and filtered coefficients}

For a topological group \(G\) and a continuous coefficient module \(M\),
we use the inhomogeneous continuous cochain complex
\[
 C^q_{\hthreects}(G,M)=C_{\hthreects}(G^q,M)
\]
with its usual differential. Complete power-series and valuation
lattices have their linear topologies. Their discrete quotients have
the quotient topology. For a profinite parameter space \(Q\), all
constructions are also made with coefficients
\(C_{\hthreects}(Q,M)\).

A filtered localization or algebraic union is interpreted at the level
of complexes. If \(M=\colim_\lambda M_\lambda\) is such a coefficient
object, our notation means
\begin{equation}\label{h3:eq:framework-filtered-cochains}
 C^\bullet_{\hthreects}(G,M)
 :=\colim_\lambda C^\bullet_{\hthreects}(G,M_\lambda).
\end{equation}
Thus a cochain has a common finite stage or pole bound. This is the
coefficient convention produced by ordinary tensor products in the
residue tests. It is generally different from the complex of all
metric-continuous maps into a topologized union.

We use completed group-ring resolutions to compute continuous
cohomology only after comparing them with the continuous cochain model
on the coefficient category in question. For complete linearly
topologized modules the comparison, compact duality, and the required
strictness statements are established in
\Cref{h3:sec:coheight-one}. For the geometric integral complexes in
\Cref{h3:sec:rational}, profinite parameters are retained through descent
and the derived limits. The finite solid resolution argument there
justifies the later calculation on their point-valued cohomology rows.

The holomorphic coefficient ring \(B_{\hthreehol}\) consists of integral-exponent
Laurent series converging on
\[
 0<|x|<1,\qquad |y|<1
\]
over the trivially valued field \(k\). We give it the inverse-limit
topology from Gauss norms on an exhaustion by closed polyannuli.
The analogous analytic completed perfection is \(B_{\hthreepk}\).
The precise exhaustion, coefficient growth conditions and root
projections are specified in \Cref{h3:sec:ordinary}. Bounded-pole
submodules of \(B_0\) carry their compact lattice topology. We
distinguish the coheight-one Cartier operator \(\mathcal C\) from
the ordinary root-projection operator \(\mathcal P\). Their
different Frobenius identities will be proved on their respective
coefficient modules.

\subsection{Geometric coefficient theory}

We use perfectoid tilting and the Fargues--Fontaine curve over an
algebraically closed perfectoid field. The notation
\(\mathcal O^\flat\) denotes the analytic tilted structure sheaf;
an integral valuation lattice is indicated separately. In particular,
the affinoid perfectoid acyclicity used below concerns this analytic
sheaf. We do not replace an integral almost-vanishing statement by
ordinary vanishing.

Let
\[
 V_s=\mathcal O(1/s)
\]
be the stable rank-\(s\), degree-one bundle. We use the classification
of vector bundles by slopes, \(H^0(\mathcal O)=\Q_p\), and the
relative vanishing of \(H^1\) for trivial bundles and positive
basic bundles as a \(v\)-sheaf. On a fixed basic stratum, bundles admit \(v\)-local
standard forms with the indicated locally profinite automorphism
groups. Slope-zero bundles correspond to rational local systems;
their integral lattices are additional data. These are the precise
background bundle results needed from
\cite{FarguesFontaine,KedlayaLiu,ScholzeWeinstein}; the precise relative arguments appear in
\Cref{h3:geom:curve-inputs}.

The Honda universal cover identifies the section sheaf of \(V_s\)
with the perfected open Honda ball, compatibly with addition and
endomorphisms. The Banach--Colmez description is provided by
\cite{LeBras,ScholzeWeinstein}. For maps over a varying perfectoid base we use
relative full faithfulness of derived sections
\cite[Corollary~3.11]{AnschuetzLeBrasFourier}; its application to positive bundles
is made explicit in \Cref{h3:sec:towers}.
\Cref{h3:sec:towers} verifies the finite-level coordinate comparison
and the relative integral frame reductions used in this paper.
\Cref{h3:sec:analytic} proves the nonproper relative support calculations;
they are not supplied by affinoid acyclicity alone.

\subsection{The two boundary calculations}

We use two boundary calculations. The rational height-three
calculation is
\begin{equation}\label{h3:eq:framework-rational-boundary}
 \pi_*L_0T
 \simeq \Lambda_{\Q_p}(\alpha_1,\alpha_3,\alpha_5),
 \qquad |\alpha_i|=-i.
\end{equation}
For the primes \(p\geq5\) considered here, \Cref{h3:prop:rational-boundary}
gives the argument of \cite[Theorem~A]{BSSW}, including its arithmetic
and tower-comparison inputs. Its identification with the particular
constant trace class needed here is treated separately in
\Cref{h3:sec:rational}; a rational dimension calculation cannot identify
a localization map.

For context, the known odd-prime height-two splitting is
\begin{equation}\label{h3:eq:framework-height-two-boundary}
 L_1R_2\simeq
 L_1(S\vee\Sigma^{-1}S)
 \vee L_0(\Sigma^{-3}S\vee\Sigma^{-4}S),
\end{equation}
as stated in \cite[Theorem~6.7]{BB}.  Exactness also gives
\(R_2/p\simeq L_{K(2)}(S/p)\), where \(S/p\) is the finite Moore
spectrum.  For \(p\geq5\), \cite[Theorem~15.1]{HS99} therefore implies
that \(\pi_a(R_2/p)\) is finite for every integer \(a\).
The finite groups \(\pi_a(R_2/p)\) give one boundary for the
intermediate finiteness argument in \Cref{h3:sec:ordinary}.  That
argument derives the required finiteness of \(L_1(R_2/p)\) directly
from the height-two group calculation of \cite[Remark~7.8]{Behrens2012}.
The coheight-one test, the ordinary comparison, and all the attaching
maps are proved below.

Morava descent will be used with an explicit convergence statement:
the completed Amitsur error tower has a uniform nilpotence bound, so
the exact tests in this proof preserve its totalization, as derived
from descendability in \Cref{h3:prop:exact-descent}.  The cooperation
calculation of Devinatz--Hopkins
\cite[Theorem~2(ii) and Proposition~2.2]{DevinatzHopkins} is used
first for the usual coefficient ring $W(\F_{p^n})$ and extended
stabilizer.  The passage to the chosen $k$, the relative algebra over
$S_k$, and its ordinary-stabilizer action-cobar terms are local
arguments of \Cref{h3:sec:descent}, which also retains the actual unit
and cofaces.
